\section*{Introduction}

Ce rapport s'inscrit dans le cadre du projet \textit{Data Protection} pour l'année académique 2024-2025. L'objectif principal de ce projet est d'appliquer une chaîne d'analyse de données sur un ensemble de données réseau ert de données physiques.

Le projet vise à explorer, préparer et évaluer ces données afin de développer des modèles de détection d'anomalies performants pour identifier divers types d'attaques sur le système étudié. Il s'appuie sur l'application et la comparaison de plusieurs algorithmes d'apprentissage supervisé, notamment \texttt{KNN}, \texttt{CART}, \texttt{Random Forest}, \texttt{XGBoost} et \texttt{MLP}. Les performances de ces modèles ont été analysées selon des métriques adaptées aux données équilibrées et déséquilibrées, tout en considérant les contraintes liées aux ressources de l'ordinateur.